\section{The method}\label{sec:method}

We never deform a known polygon. Instead we build the polygon from its edge
lengths. Write $2h_i$ for the length of the edge from vertex $i$ to vertex
$i+1$. A chord of length $2h$ in a circle of curvature $\kappa$ makes the angle
$\arcsin(\kappa h)$ with the tangent at either end, so if the circle through
three consecutive vertices is to have curvature $\kappa_i$, the polygon must
turn at vertex $i$ by
$\theta_i=\arcsin(\kappa_ih_{i-1})+\arcsin(\kappa_ih_i)$. Laying the edges
end to end with these turns, the polygon closes up after one revolution
exactly when three real equations hold: the two coordinates of the closure
gap and the total turn $\sum\theta_i=2\pi$. It therefore suffices to solve
three equations in the $n$ positive half-lengths $h_i$, subject to the bounds
$\kappa_i h_{i-1}<1$ and $\kappa_i h_i<1$. Section~\ref{sec:setup} proves
that every such solution gives a simple, strictly convex, coherent polygon.

\begin{figure}[t]
\centering
\figfile{method}{\textwidth}
\caption{The method, computed for the word
$\mu=(1,\allowbreak\,0.94,\allowbreak\,0.63,\allowbreak\,0.44,\allowbreak\,1.75,\allowbreak\,1.88,\allowbreak\,1.25,\allowbreak\,0.56,\allowbreak\,0.5)$.
(a) The frame of four active edges $e_0,\dots,e_3$ (shaded) and their eight
endpoint readings (black); the other readings (grey) are dormant.
(b) A regular anchor: the dormant edges have length zero, each block of
vertices between two active edges is a single point (open circles), and the
polygon is a quadrilateral. (c) The implicit function theorem opens the
dormant edges to small positive lengths, keeping every prescribed curvature.
(d) Another realization of the same word, for comparison.}
\label{fig:method}
\end{figure}

The main idea is to put all the length on four edges
(Figure~\ref{fig:method}). If every other edge has length zero, the vertices
between two consecutive long edges collapse to one point. The turn at such a
collapsed block involves only the curvatures at its two ends, and the three
equations become three equations in four unknowns that involve only the
eight curvature readings at the ends of the four chosen edges. A positive
solution at which the three equations have full rank is a \emph{regular
anchor}. It is a closed quadrilateral whose corners carry two circles each.
The implicit function theorem then opens every edge of length zero to a small
positive length while keeping all the prescribed curvatures, and the result
is a realization (Section~\ref{sec:anchor}). The same argument, with the
curvatures as parameters, gives the local $C^1$ dependence.

It remains to choose the four edges so that a regular anchor exists. After
removing repeated consecutive curvatures and rescaling, this is a
combinatorial problem about the word $\kappa_1,\dots,\kappa_n$, solved by one
of three explicit constructions, each on a frame of the shape
$(0,\allowbreak\,i-1,\allowbreak\,i,\allowbreak\,n-1)$: the two edges at a peak with curvature $1$ at vertex
$0$ and the two edges at a vertex $i$ further on. The resulting four-length
system depends on five numbers only (Section~\ref{sec:system}).
\begin{description}
\item[(R)] If the maximum of $\kappa$ is attained twice, take the edges at
  the two maxima. An explicit quadrilateral with perpendicular diagonals,
  deformed in two parameters, gives the anchor.
\item[(F)] Otherwise let $0$ be the highest peak other than the global
  maximum $H$. A number $\varphi$ built from the readings next to $0$ and $H$,
  the \emph{forward flag}, decides the construction. If it is nonnegative,
  take the edges at the \emph{first transition}: the first vertex $i$ on the
  way up to $H$ at which the flag of $(\kappa_{i},\kappa_{i+1})$ becomes
  nonnegative.
\item[(N)] If the forward flag and its mirror image, the backward flag, are
  both negative, take the edges at the two peaks $0$ and $H$ and again an
  orthogonal quadrilateral.
\end{description}
Each construction finds a zero by the intermediate value theorem or the
Poincar\'e--Miranda theorem, and proves the rank condition by an explicit
computation (Section~\ref{sec:constructions}).
